% ---- 字体：正文思源宋体（Noto Serif SC），代码 Cascadia Code ----
\usepackage{fontspec}
% 思源/Noto 在本机是可变字体(VF)，按「族名」加载会让 xdvipdfmx 崩(Invalid font -1)，
% 必须按「文件名」加载。
% 西文正文用思源宋体自带的拉丁衬线，与中文同源、观感统一。
\setmainfont{NotoSerifSC-VF.ttf}[AutoFakeSlant=0.15]
% 中文正文：思源宋体
\setCJKmainfont{NotoSerifSC-VF.ttf}[AutoFakeSlant=0.15]
% 无衬线：思源黑体
\setCJKsansfont{NotoSansSC-VF.ttf}
% 等宽（代码）：Cascadia Code，符号覆盖广（∈ ≠ ≡ ⌊ ⌋ ∘ 等），顺带消掉缺字警告。
% 【务必按文件名从 TeX Live 静态 OTF 加载】：本机同时装了 Windows 版可变字体
% CascadiaCode.ttf，按族名 {Cascadia Code} 会命中偏粗默认实例。Path 必须 Windows
% 盘符风格（xelatex 原生 Windows 程序，读不懂 MSYS 的 /d/...）。
% 【字重取 SemiLight 且把 Bold 也映射到 SemiLight】：Cascadia 的 Regular 相对思源宋体
% 偏重，行内代码 `vector` 挨着正文时会显得“加粗”；标题是 \bfseries，若 Bold face 更粗，
% 标题里的行内代码会真加粗并顺带污染目录。统一用 SemiLight 且不给独立 Bold face，
% 让行内代码在正文/标题/目录里都与宋体粗细相称、任何语境都不再变粗。
% （jblight 主题仅 Alert 一个 token 用到 bold，C++ 代码里几乎不出现，无损。）
\setmonofont{CascadiaCode}[
  Path=D:/texlive/2026/texmf-dist/fonts/opentype/public/cascadia-code/,
  Extension=.otf,
  UprightFont=*-SemiLight,
  BoldFont=*-SemiLight,
  ItalicFont=*-SemiLightItalic,
  BoldItalicFont=*-SemiLightItalic,
  Scale=0.90]
% 代码块里的中文注释回退到思源宋体
\setCJKmonofont{NotoSerifSC-VF.ttf}[Scale=0.95]

% ---- 代码块：长行自动折断、缩小字号；中文注释回退到主 CJK 字体正常显示 ----
\usepackage{fvextra}
\fvset{breaklines=true,breakanywhere=true,fontsize=\small,numbersep=8pt}

% ---- 代码块外观：圆角边框 + 行号排在框内左侧 gutter、竖线分隔（对标 IDE / 截图） ----
% pandoc 生成的代码块外层是 Shaded 环境（默认 framed 的 snugshade）。header-includes
% 在 pandoc 的高亮宏之后插入，故此处 \renewenvironment 覆盖生效，shadecolor 也已由
% pandoc 按主题背景色定义好，直接拿来做底色。
\usepackage{tcolorbox}
\tcbuselibrary{breakable,skins}
\definecolor{codeframe}{HTML}{D9DDE3}
\definecolor{coderule}{HTML}{CBD0D7}
\definecolor{codenum}{HTML}{9AA0A8}
% 行号：浅灰、小一号
\renewcommand{\theFancyVerbLine}{\color{codenum}\footnotesize\arabic{FancyVerbLine}}
% gutter 宽度：容纳行号 + 与代码的间距
\newlength{\codegutter}\setlength{\codegutter}{2.7em}
\renewenvironment{Shaded}{%
  \begin{tcolorbox}[breakable, enhanced, colback=shadecolor, colframe=codeframe,
    boxrule=0.4pt, arc=3pt, outer arc=3pt, boxsep=0pt,
    left=\codegutter, right=6pt, top=4pt, bottom=4pt,
    underlay unbroken and first={\draw[coderule, line width=0.5pt]
      ([xshift=2.05em]interior.north west) -- ([xshift=2.05em]interior.south west);},
    underlay middle and last={\draw[coderule, line width=0.5pt]
      ([xshift=2.05em]interior.north west) -- ([xshift=2.05em]interior.south west);}]%
}{\end{tcolorbox}}

% ---- 段落用间距代替首行缩进，贴近 WIDA 的紧凑观感 ----
\usepackage{parskip}

% 页眉只显示当前一级章节；进入新章时清空旧的右标记。
\usepackage{fancyhdr}
\pagestyle{fancy}
\fancyhf{}
\renewcommand{\sectionmark}[1]{\markboth{#1}{}}
\renewcommand{\subsectionmark}[1]{}
\fancyhead[L]{\nouppercase{\leftmark}}
\fancyhead[R]{\thepage}
\renewcommand{\headrulewidth}{0.4pt}
\setlength{\headheight}{14pt}
\usepackage{etoolbox}
\pretocmd{\section}{\clearpage}{}{}

% 目录已超过一千页，默认页码盒只够三位数；扩宽后四位页码不会越过右边界。
\makeatletter
\renewcommand{\@pnumwidth}{2.4em}
\renewcommand{\@tocrmarg}{3.2em}
\makeatother

% ---- 拉丁扩展字符兜底：思源宋体缺匈牙利语 ő（Erdős-Gallai 定理里的人名），
%      正文（非代码）里遇到就回退到 Latin Modern（TeX Live 必带，覆盖全）。----
\usepackage{newunicodechar}
\newfontfamily\latinfallback{Latin Modern Roman}
\newunicodechar{ő}{{\latinfallback ő}}
\newunicodechar{Ő}{{\latinfallback Ő}}
